% =====================================================================
%   附录 E：数字卡牌游戏 (Digital Card Game)
% =====================================================================

\section{数字卡牌游戏}\label{app:dcg}

\subsection{数据集细节}

\subsubsection*{构造细节}

本文以 Aquawar 框架作为交互系统的基础。第一类交互为动作阶段：模型需选择
其要执行动作的鱼，然后选取技能的目标。为保证模型操作的有效性，我们对合法
动作进行检查。第二类交互为猜测阶段：我们向模型提供已知信息，包括鱼的种类
与技能描述、敌方目标等。我们为测试目的提供了两种朴素策略。游戏流程的
详细定义与说明如下：

\begin{itemize}
    \item \textbf{玩家与卡牌}：这是一款双人对战游戏，每队拥有 4 条宠物鱼
          （即卡牌）。牌池共包含 10 种鱼（详见附录~\ref{app:dcg}~E.2），
          双方玩家在游戏开始前需各自从牌池中挑选 4 条确定的鱼。
    \item \textbf{初始状态}：每条鱼拥有 400 点初始生命值、200 点初始攻击力、
          一项主动技能与一项被动技能。
    \item \textbf{基本规则}：玩家每回合选择一条存活的鱼，使用其主动技能或
          普通攻击敌方一条鱼。任何存活鱼的被动技能在满足触发条件时会自动
          触发。
    \item \textbf{断言机制}：玩家鱼的种类身份在最初是隐藏的。对手玩家可在
          每一回合猜测玩家一条鱼的身份；若猜测正确，则该鱼的身份被公开，
          且玩家所有鱼均会受到伤害。
    \item \textbf{回合流程}：在一个回合内，该回合的玩家首先断言敌方一条
          存活且身份未公开的鱼的身份；若断言正确，则敌方所有存活的鱼
          均受到伤害。随后，该玩家可以命令一条存活的鱼执行一次普通攻击或
          主动技能。此后，任何满足触发条件的鱼将释放其被动技能。
    \item \textbf{胜利条件}：游戏结束时，存活鱼数量较多的一方获胜。
\end{itemize}

为兼顾智能体的参与度与游戏的复杂度，我们设计了两阶段的游戏逻辑：第一阶段
移除了断言机制，第二阶段则保留断言机制。我们分别在两个阶段上对所有模型
进行测试，并取两阶段的平均表现作为最终得分。

我们选择两种朴素的玩法策略作为基线：
\begin{itemize}
    \item \textbf{策略一}：从所有可选动作空间中简单地随机选择动作。
    \item \textbf{策略二}：在可能的情况下优先尝试 AOE 攻击，并持续评估能否
          实现一击必杀；接着尝试使用主动技能，最后退而采用普通攻击。总体
          上，该策略遵循一定的模式，但未必是最优策略。
\end{itemize}

\subsubsection*{评测设置}

每次游戏对局中，我们按以下步骤进行评估：
\begin{enumerate}
    \item \textbf{初始化}：我们在 Ubuntu 20.04 环境下初始化经过修改的游戏
          逻辑环境（基于 pybind 编译）以及基线游戏智能体。
    \item \textbf{交互}：我们根据不同游戏阶段在指令提示中放置规则描述，
          LLM 智能体在游戏逻辑环境中与基线进行策略性对抗。我们给予 LLM
          智能体 5 次按正确格式响应的机会；若在该次数内未能输出合法动作，
          则立即判定其失败。同时，我们鼓励模型以 CoT 形式输出其推理过程。
    \item \textbf{结果计算}：在交互过程中，我们将记录完整的游戏过程以供
          对战回放，并计算游戏结果以获得任务的评测指标。
\end{enumerate}

\subsubsection*{评测指标}

我们采用的评测指标涵盖基本玩法要素，包括获胜回合数（Win Round）、总回合数
（Total Round）、胜率（Win Rate）、总伤害占总生命值的比例（Damage Rate），
最终根据以上指标计算综合奖励分数：
\begin{equation}\label{eq:reward}
    \text{reward} = 0.7 \times \text{metric}_\text{winrate} + 0.3 \times \text{metric}_\text{damagerate}.
\end{equation}

\subsection{鱼类属性}

依据游戏规则，本游戏共包含以下 10 种鱼：

\begin{itemize}
    \item \textbf{Spray（水柱喷射鱼）}
        \begin{itemize}
            \item Counter（被动）：当一名队友的生命值低于 30\% 时，对攻击者
                  造成 30 点伤害。
            \item AOE（主动）：对其攻击力的 35\% 对所有敌人造成伤害。
        \end{itemize}
    \item \textbf{Flame（烈焰鱼）}
        \begin{itemize}
            \item Counter（被动）：当一名队友的生命值低于 30\% 时，对攻击者
                  造成 30 点伤害。
            \item Infight（主动）：对一名存活的队友造成 75 点伤害，并使自身
                  攻击力提升 140 点。
        \end{itemize}
    \item \textbf{Eel（鳗鱼）}
        \begin{itemize}
            \item Deflect（被动）：受到攻击时，将 70\% 的伤害转移给队友，
                  自身仅承受 30\% 的伤害；在累计承受 200 点伤害后，获得 40 点
                  攻击力加成。
            \item AOE（主动）：对其攻击力的 35\% 对所有敌人造成伤害。
        \end{itemize}
    \item \textbf{Sunfish（翻车鱼）}
        \begin{itemize}
            \item Deflect（被动）：受到攻击时，将 70\% 的伤害转移给队友，
                  自身仅承受 30\% 的伤害；在累计承受 200 点伤害后，获得 40 点
                  攻击力加成。
            \item Infight（主动）：对一名存活的队友造成 75 点伤害，并使自身
                  攻击力提升 140 点。
        \end{itemize}
    \item \textbf{Barracuda（梭鱼）}
        \begin{itemize}
            \item Reduce（被动）：每次受到攻击时，有 30\% 的概率完全规避该
                  伤害。
            \item Crit（主动）：对一名敌人造成 120 点暴击伤害。
        \end{itemize}
    \item \textbf{Mobula（蝠鲼）}
        \begin{itemize}
            \item Reduce（被动）：每次受到攻击时，有 30\% 的概率完全规避该
                  伤害。
            \item Subtle（主动）：选择一名队友或自身，使其在受到攻击时所受
                  伤害减少 70\%，并使其攻击力提升 20 点。
        \end{itemize}
    \item \textbf{Octopus（章鱼）}
        \begin{itemize}
            \item Heal（被动）：当受到攻击且生命值仍大于 0 时，恢复 20 点
                  生命值。
            \item Infight（主动）：对一名存活的队友造成 75 点伤害，并使自身
                  攻击力提升 140 点。
        \end{itemize}
    \item \textbf{Whiteshark（白鲨）}
        \begin{itemize}
            \item Heal（被动）：当受到攻击且生命值仍大于 0 时，恢复 20 点
                  生命值。
            \item Crit（主动）：对生命值最低的敌人造成相当于自身攻击力 120\%
                  的暴击伤害；若目标生命值低于 160，则暴击伤害提升至 140\%。
        \end{itemize}
    \item \textbf{Hammerhead（锤头鲨）}
        \begin{itemize}
            \item Explode（被动）：当受到攻击且未死亡时，对攻击来源造成 40
                  点伤害；当生命值低于 20\% 时，攻击力额外提升 15 点。
            \item Crit（主动）：对生命值最低的敌人造成相当于自身攻击力 120\%
                  的暴击伤害；若目标生命值低于 160，则暴击伤害提升至 140\%。
        \end{itemize}
\end{itemize}

可以看出，不同宠物鱼的主动与被动技能之间存在重叠。这一设计旨在更好地隐藏
宠物鱼的身份信息，并增强游戏的策略深度。

\subsection{提示词示例}

我们使用如下格式的提示词来执行动作：
\begin{quote}\small\ttfamily
This is a two-player battle game with four pet fish on each team. The
types of fish may vary.\\
Each fish has its 400 initial health, 200 attack power, active ability,
and passive ability. You can choose a live fish to use its active skill or
normal attack (causing half of attack power as damage) on an enemy fish
each round. When the conditions are met, the fish's passive ability will
automatically trigger, regardless of whether it is chosen.\\
Your fish's identity is initially hidden. The enemy can guess one of your
fish's identity in each round. If the enemy guesses right, your fish's
identity is revealed, and each of your fish will get 50 damage.\\
The victory condition is to have more fish alive at the end of the game.\\
\\
The following are the four types of your pet fish: \{\dots\}\\
The following are the four types of enemy's pet fish: \{\dots\}\\
\\
Play the game with me. In each round, you should output your thinking
process, and return your move with following JSON format:\\
\verb|{'pick_fish': 'pick an alive fish, you should give the name of the|
\verb|alive fish', 'action': 'choose from [normal, active]',|
\verb|'target_position': "target's position, you must choose from [0,3]"}|\\
\\
Notice! You must return your move in each round. Otherwise, you will be
considered defeated.
\end{quote}

第二阶段的断言提示词与上述类似，但将动作格式替换为对敌方鱼身份的猜测。

\subsection{对战生成}

在本文正文的研究中，我们为公平评测起见，固定了若干预设，使得所有模型均
使用相同的对战场景进行测试。然而，这些任务的生成实际上可引入一定的
随机性。对战中的随机性主要体现在以下方面：
\begin{itemize}
    \item \textbf{队伍构成}：双方宠物鱼的选择，包括其数量与种类。
    \item \textbf{属性确定}：双方各宠物鱼的具体属性数值。
\end{itemize}

我们允许选择不同难度等级以生成随机对战。本节将进一步解释"难度"这一概念的
运作方式。

\subsubsection*{战斗力（Combat Power）}

为有效评估战斗力，关键是把握其基本原理：在两支队伍的对抗中，战斗力较强
的一方通常具有更高的获胜概率。

记鱼 $c$ 的生命值为 $\text{HP}_c$，其攻击力（即普通攻击的期望伤害）为
$\text{ATK}_c$。设想一个简单场景：来自我方队伍的一条鱼 $f$ 与来自敌方队伍的
一条鱼 $h$ 互相对战，同时对彼此造成伤害。双方能够承受战斗的持续时间分别
为 $\text{HP}_f / \text{ATK}_h$ 与 $\text{HP}_h / \text{ATK}_f$。
在此情形下，若 $\text{HP}_f / \text{ATK}_h \ge \text{HP}_h / \text{ATK}_f$，
即 $\text{HP}_f \cdot \text{ATK}_f \ge \text{HP}_h \cdot \text{ATK}_h$，
我方队伍更有可能获胜。

因此，对于单鱼队伍，我们将其战斗力定义为 $\text{HP} \cdot \text{ATK}$，这与
上述判据相符。

考虑由多条鱼组成的队伍 $T$，在排除任何特殊技能的情况下，每一回合需从我方
队伍中选择一条鱼攻击敌方的一条鱼。由此可将整支队伍视为整体，总生命值为
所有鱼生命值之和，平均攻击力（在正常情况下，应选择攻击力最高的鱼，但
考虑到对手通常会优先攻击攻击力最高的鱼以削弱其持续作战能力，我们假设
所有鱼的攻击频次相当）则被纳入考量。

由此，队伍 $T$ 的战斗力 $\text{Power}(T)$ 定义为：
\begin{equation}\label{eq:power}
    \text{Power}(T) = \frac{1}{\#T} \sum_{c \in T} \text{HP}_c \sum_{c \in T} \text{ATK}_c.
\end{equation}

\subsubsection*{难度：战斗力之比}

对战难度通过对战双方队伍的战斗力之比来量化：
\begin{equation}\label{eq:rho}
    \rho(H \mid F) = \frac{\text{Power}(H)}{\text{Power}(F)}.
\end{equation}

式中，$F$ 代表友方队伍，$H$ 代表敌方队伍。可见，比值 1 表示难度均衡，即
双方实力对等；进一步地，该难度指标意味着 $\rho(T_1 \mid T_2)$ 个 $T_2$ 队伍
在实力上等同于 1 个 $T_1$ 队伍。

借助这一框架，研究者能够为评估各种策略的胜率精确设定具体难度等级，从而
有效衡量其效能。
